\documentclass[12pt]{article}
\usepackage[utf8]{inputenc}
\usepackage{amsmath, amssymb}
\usepackage{enumitem}
\usepackage{fancyhdr}
\usepackage{geometry}
\geometry{a4paper, margin=1in}
\pagestyle{fancy}
\fancyhf{}
\rhead{Toán Rời Rạc - Phần Logic}
\lhead{Đại học Phenikaa}
\cfoot{\thepage}

\title{Giải Bài Tập Toán Rời Rạc \newline Phần 1: Logic}
\date{Năm học 2024 -- 2025}

\begin{document}

\maketitle

\section*{Bài 1.1. Chứng minh các biểu thức sau là hằng đúng}

\begin{enumerate}[label=\alph*)]
    \item \(((P \rightarrow Q) \land P) \rightarrow Q\)
    \\ \textbf{Cách 1: Bảng chân trị}\\
    \begin{tabular}{|c|c|c|c|c|}
        \hline
        P & Q & $P \rightarrow Q$ & $(P \rightarrow Q) \land P$ & Toàn mệnh đề \\
        \hline
        T & T & T & T & T \\
        T & F & F & F & T \\
        F & T & T & F & T \\
        F & F & T & F & T \\
        \hline
    \end{tabular}\\
    \textbf{Cách 2:} Theo luật Modus Ponens.

    \item $P \land Q \rightarrow P$\\
    \begin{tabular}{|c|c|c|c|}
        \hline
        P & Q & $P \land Q$ & Mệnh đề \\
        \hline
        T & T & T & T \\
        T & F & F & T \\
        F & T & F & T \\
        F & F & F & T \\
        \hline
    \end{tabular}

    \item $\neg(P \land Q) \land P \rightarrow \neg Q$\\
    Biến đổi: $\neg(P \land Q) \equiv \neg P \lor \neg Q$, mệnh đề luôn đúng theo bảng chân trị.

    \item $(P \rightarrow Q \land R) \rightarrow ((P \rightarrow Q) \land (P \rightarrow R))$\\
    Theo định nghĩa kéo theo, mệnh đề là hằng đúng.

    \item $((P \land Q) \leftrightarrow P) \rightarrow (P \rightarrow Q)$\\
    \begin{tabular}{|c|c|c|c|c|}
        \hline
        P & Q & $P \land Q$ & $\leftrightarrow$ & Toàn mệnh đề \\
        \hline
        T & T & T & T & T \\
        T & F & F & F & T \\
        F & T & F & T & T \\
        F & F & F & T & T \\
        \hline
    \end{tabular}
\end{enumerate}

\section*{Bài 1.2. Quy tắc suy diễn}
\begin{enumerate}[label=\alph*)]
    \item $((X_1 \rightarrow X_2) \land (\neg X_3 \lor X_4) \land (X_1 \lor X_3)) \rightarrow (\neg X_2 \rightarrow X_4)$\\
    Dùng các luật logic và suy diễn như Modus Tollens để biến đổi và chứng minh mệnh đề là hằng đúng.

    \item \textbf{Văn bản:}\\
    \begin{itemize}
        \item $A$: An được thưởng cuối năm
        \item $B$: An đi Đà Lạt
        \item $C$: An thăm Thiền viện
        \item Giả thiết: $A \rightarrow B$, $B \rightarrow C$, $\neg C$\\
        \textbf{Kết luận:} $\neg A$\\
        Suy luận đúng nhờ phép phản chứng (Modus Tollens).
    \end{itemize}
\end{enumerate}

\section*{Bài 1.3. Dịch các câu thành biểu thức logic}
\begin{enumerate}[label=\alph*)]
    \item $\forall x (\text{Chim ruồi}(x) \rightarrow \text{Sặc sỡ}(x))$
    \item $\neg \exists x (\text{Chim lớn}(x) \land \text{Ăn mật}(x))$
    \item $\forall x ((\text{Chim lớn}(x) \land \neg \text{Ăn mật}(x)) \rightarrow \text{Màu xám}(x))$
    \item $\forall x (\text{Chim ruồi}(x) \rightarrow \text{Nhỏ}(x))$
\end{enumerate}

\section*{Bài 1.4. Dịch bằng lượng từ với $L(x, y)$ là “x yêu y”}
\begin{enumerate}[label=\alph*)]
    \item $\forall x\, L(x, \text{Mai})$
    \item $\forall x\, \exists y\, L(x, y)$
    \item $\exists x\, \forall y\, L(y, x)$
    \item $\neg \exists x\, \forall y\, L(x, y)$
    \item $\exists x (\forall y\, \neg L(x, y) \lor \forall y\, \neg L(y, x))$
    \item $\exists x\, \neg L(\text{Nam}, x)$
    \item $\exists! x\, \forall y\, L(y, x)$
    \item $\exists x \exists y (x \neq y \land L(\text{Tuấn}, x) \land L(\text{Tuấn}, y) \land \forall z\, (L(\text{Tuấn}, z) \rightarrow (z = x \lor z = y)))$
\end{enumerate}
